\documentclass[12pt,a4paper,twoside]{article}

% are used to comment out text

%------------------------------------------------------------------------------------------------------------
%EDIT THESE FIELDS WITH YOUR DETAILS
%------------------------------------------------------------------------------------------------------------
\newcommand{\studentname}{Chen Zhi Shen}
\newcommand{\studentnumber}{1644420}

\title{Lab 3 Report}

%\newcommand{\tutor}{Fan Hao}
\newcommand{\labgroup}{M1.03}
\newcommand{\labdate}{28th October 2024}
\newcommand{\reportdate}{\today}
%------------------------------------------------------------------------------------------------------------

%------------------------------------------------------------------------------------------------------------
%PACKAGES. PACKAGES ARE USED TO EXTEND THE FUNCTIONALITY OF LATEX. YOU SHOULDN'T
%NEED ANY MORE, SPEND YOUR TIME ON THE CONTENT!
%------------------------------------------------------------------------------------------------------------

\usepackage[a4paper,left=25mm,right=25mm,bottom=25mm,top=25mm]{geometry}
\usepackage{fancyhdr}
\usepackage{subfigure}
\usepackage{graphicx}
\usepackage{cleveref}
\usepackage{booktabs}
\usepackage[british]{babel}
\usepackage[square,comma,numbers,sort&compress]{natbib}
\usepackage{csvsimple}
\usepackage{graphicx}
\usepackage{pgfplotstable,filecontents}
\usepackage{gensymb}
\usepackage{array}
\usepackage{tabu}
\usepackage{multirow}
\usepackage{url}
\pgfplotsset{compat=1.9}% supress warning
\begin{document}

%------------------------------------------------------------------------------------------------------------

%------------------------------------------------------------------------------------------------------------
%DO NOT EDIT THESE FIELDS
%------------------------------------------------------------------------------------------------------------
\setlength{\parindent}{0em}
\setlength{\parskip}{0.5em}
\author{\studentname ~ \studentnumber. Lab group: \labgroup.}
\date{Lab date: \labdate. Report date: \reportdate.}

\renewcommand\abstractname{Summary}

\pagestyle{fancy}
\fancyhead{}
\fancyhead[C]{\studentname~\studentnumber}
\fancyfoot{}
\fancyfoot[C]{\thepage}
\renewcommand{\headrulewidth}{0pt}
\renewcommand{\footrulewidth}{0pt}

\renewcommand\bibname{References}

%------------------------------------------------------------------------------------------------------------

%\renewcommand\abstractname{Abstract}
\maketitle

%------------------------------------------------------------------------------------------------------------
%YOUR REPORT BEGINS HERE
%------------------------------------------------------------------------------------------------------------
\section{Lab Introduction and Team Information}
The aim of this laboratory work\cite{Paper2014} is to ...

Our team members include ...

Tasks for team members.

\section{Individual Work}
Question 1:

Answer 1:

Question 2:

Answer 2:

\section{Team Work}
Task 1:

Answer 1:

Task 2:

Answer 2:

\section{Discussion}
...

\section{Conclusions}
...
%------------------------------------------------------------------------------------------------------------
%YOUR REPORT ENDS HERE, THEN THE BIBLIOGRAPHY IS CREATED AND DOCUMENT ENDED
%------------------------------------------------------------------------------------------------------------

\bibliographystyle{unsrtnat}
\bibliography{refs}

%------------------------------------------------------------------------------------------------------------
%BELOW ARE JUST EXAMPLE SECTIONS TO HELP YOU - PLEASE DELETE THEM
%------------------------------------------------------------------------------------------------------------

\end{document}
